% 02_introduction.tex — Introduction
% Drafted Session 4bq — block-and-wait review gate.
% **DRAFT v1 — awaiting owner review.**

\section{Introduction}\label{sec:introduction}

Amorphous heavy-metal alloys are of growing interest in spintronics, where
their ability to generate spin--orbit torque from a structural, rather than a
chemical, order makes them attractive spin-current sources. Ion-implanted
amorphous Pt(P), for example, exhibits a giant spin Hall effect and
spin--orbit torque that can exceed those of crystalline heavy metals
\cite{shashank2025}, and amorphous materials more broadly are emerging across
electronics applications as their properties are tuned by disorder rather
than by composition alone \cite{shi2025,madanchi2024}. Because these
properties depend on the atomic-scale structure, knowing {\em when} a
candidate system becomes amorphous --- and over what composition range --- is
central to designing such films.

Predicting structure across that regime is, however, costly. Generating an
amorphous configuration is conventionally done by first-principles melt-quench
molecular dynamics, which is prohibitively expensive at the level of accuracy
needed for a multi-element alloy \cite{meltquench2026,madanchi2024}, and the
tens of thousands of reference calculations that a general machine-learned
interatomic potential demands are out of reach for exploratory screening
\cite{activemlp2020}. Evolutionary and machine-learning-assisted sampling
have lowered the barrier \cite{artrith2018,nahas2016,valladares2011}, but a
practical route to ask ``does this interstitial alloy amorphize, and at what
loading?'' remains open.

Here we take a reference-anchored approach. Rather than melting and quenching,
we model the amorphous state as the low-energy basin around a chosen
crystalline reference and enumerate that basin with a surrogate-guided
global-optimization search
\cite{gofee2020,gofee2022,agox2022}. The crystalline reference is a physical
choice the user controls --- unique for a host with a single stable phase, and
a genuine design choice for a polymorphic one. Amorphization is then confirmed
structurally with the Steinhardt $q_6$ bond-orientational order parameter
\cite{steinhardt1983}, cross-checked against simulated X-ray diffraction.
Because the surrogate guides the search, the basin is sampled at a small
fraction of the cost of melt-quench dynamics \cite{artrith2018}.

Applied to interstitial platinum phosphide (Pt(P)), whose fcc host has a
single stable phase, the screening locates the onset of structural
amorphization at 20--30~at.\\% P, saturates with further loading, and is
mildly promoted by cell expansion --- a prediction consistent with the
experimental amorphization onset of ion-implanted Pt(P) \cite{shashank2025}.
The reference-anchored screening thus offers a cheap, structurally verified
route to the onset of amorphization in interstitial alloys.
